\documentclass[11pt]{article}
\usepackage[a4paper,margin=28mm]{geometry}
\usepackage{amsmath,amssymb,amsthm,mathtools}
\usepackage{enumitem}
\usepackage{hyperref}
\usepackage{microtype}

\newtheorem{theorem}{Theorem}[section]
\newtheorem{lemma}[theorem]{Lemma}
\newtheorem{proposition}[theorem]{Proposition}
\newtheorem{corollary}[theorem]{Corollary}
\newtheorem{definition}[theorem]{Definition}
\newtheorem{counterexample}[theorem]{Counterexample}
\newtheorem{remark}[theorem]{Remark}

\newcommand{\Q}{\mathbb Q}
\newcommand{\Z}{\mathbb Z}
\newcommand{\F}{\mathbb F}
\newcommand{\SL}{\operatorname{SL}}
\newcommand{\GL}{\operatorname{GL}}
\newcommand{\ord}{\operatorname{ord}}
\newcommand{\rad}{\operatorname{rad}}
\newcommand{\Gal}{\operatorname{Gal}}
\newcommand{\im}{\operatorname{im}}

\title{Quantitative Frey Primes, Stacky Orbicurves, and the Exact Remaining Source Theorems in a Lean Formalization of the $abc$ Conjecture}
\author{ABCConjecture formalization project}
\date{August 2026}

\begin{document}
\maketitle

\begin{abstract}
We isolate and prove several structural and quantitative statements that occur
immediately before the unresolved anabelian, tempered, and multiradial source
theorems in an inter-universal-Teichm\"uller-theoretic route to the
$abc$ conjecture.  The principal new arithmetic route replaces a merely
eventual open-image theorem by a uniform semistable argument: Mazur
irreducibility, a Tate-inertia transvection, and an explicit elementary
factorization of determinant-one matrices imply containment of
$\SL_2(\F_\ell)$.  A two-stage Euclidean construction then selects such a
prime above an arbitrary lower threshold, avoiding a local Tate order and one
residue characteristic, with a polynomial upper bound in the local order.
Consequently the logarithm of the selected prime is logarithmic in the Frey
height.  We prove two independent absorption principles showing that either a
relative conductor estimate or a sublinear-height estimate for the genuine
IUT different/error terms is sufficient for $abc$.

On the geometric side we prove that the coarse quotient by the sign
involution cannot satisfy the required cartesian square, exhibit the missing
stacky inertia, and prove the generic action-groupoid weak-pullback theorem on
objects and morphisms.  On the multiradial side we separate the distinguished
$q$-branch needed for the lower Corollary~3.12 inequality from the all-output
upper estimate, construct a semantic transfer theorem from genuine outputs to
a generated Ind1/Ind2/Ind3 syntax, and distinguish permutation-orbit volume
from the generally larger volume of a componentwise union.  We record exact
counterexamples to coarse quotients, componentwise arbitrary-Ind2 unions, and
unrestricted positive rescaling.

Every theorem claimed as formalized is represented by a corresponding Lean
module in the repository.  The paper does not claim a complete proof of
$abc$: the remaining source theorems are stated precisely in the final
section.
\end{abstract}

\section{The logarithmic $abc$ target and the Frey curve}

Let $a,b,c\in\Z_{>0}$ satisfy
\[
  a+b=c,\qquad \gcd(a,b)=\gcd(b,c)=\gcd(c,a)=1.
\]
Write
\[
  h(a,b,c)=\log c,
  \qquad r(a,b,c)=\log\rad(abc).
\]
The logarithmic $abc$ conjecture asserts that for every $\varepsilon>0$ there
is a constant $C_\varepsilon$ such that
\[
  h(a,b,c)\le (1+\varepsilon)r(a,b,c)+C_\varepsilon.
\]
We use the integral Frey equation
\[
 E_{a,b}: y^2=x(x-a)(x+b).
\]
Its standard invariants are
\[
 c_4=16(a^2+ab+b^2),
 \qquad \Delta=16(abc)^2.
\]
The existing Lean development proves the exact rational $j$-invariant,
height corridor, discriminant--conductor comparison, and the reduction of the
public IUT IV inequality to the logarithmic $abc$ target.

\section{Odd support and the local multiplicative signature}

\begin{lemma}[Odd support outside one point]
Every primitive positive $abc$ triple either admits an odd prime
$p\mid abc$, or equals $(1,1,2)$.
\end{lemma}

\begin{proof}
If no odd prime divides any of $a,b,c$, then each is a power of two.  Pairwise
coprimality permits at most one of them to exceed one.  Since $a+b=c$ and
$a,b>0$, it follows that $a=b=1$ and $c=2$.
\end{proof}

\begin{lemma}[Odd multiplicative signature]
Let $p$ be an odd prime dividing $abc$.  Then
\[
 p\mid\Delta,
 \qquad p\nmid c_4.
\]
\end{lemma}

\begin{proof}
The first assertion follows from $\Delta=16(abc)^2$.  Suppose, for example,
that $p\mid a$.  Pairwise coprimality gives $p\nmid b$, while
\[
 a^2+ab+b^2\equiv b^2\not\equiv0\pmod p.
\]
The cases $p\mid b$ and $p\mid c=a+b$ are identical; in the latter case
$a\equiv-b\pmod p$, so
$a^2+ab+b^2\equiv a^2\not\equiv0\pmod p$.
Since $p$ is odd, the factor $16$ is invertible modulo $p$.
\end{proof}

For a DVR, an integral model with unit $c_4$ is minimal.  Indeed, under an
admissible variable change of scale $u$,
\[
 c_4' =u^{-4}c_4,
 \qquad \Delta'=u^{-12}\Delta.
\]
If the transformed model remains integral, then its $c_4$-valuation is at
most one in the multiplicative convention.  When the original $c_4$ has
valuation one, this implies $v(u^{-1})\le1$, hence
$v(\Delta')\le v(\Delta)$.  This is precisely the maximality condition in
Mathlib's definition of a minimal equation.

\begin{theorem}[Integral-model reduction criterion]
Let $R$ be a DVR with fraction field $K$ and let $W_0/R$ be a Weierstrass
equation.  Let $W/K$ be its base change.
\begin{enumerate}[label=(\roman*)]
\item If $\Delta(W_0)$ is a unit of $R$, then $W$ has good reduction.
\item If $\Delta(W_0)$ is a nonunit and $c_4(W_0)$ is a unit, then $W$ has
multiplicative reduction.
\end{enumerate}
\end{theorem}

\begin{proof}
The mapped equation is integral.  The adic valuation identities
\[
 v(r)=1\iff r\notin\mathfrak m,
 \qquad v(r)<1\iff r\in\mathfrak m
\]
translate the two ideal signatures into Mathlib's good and multiplicative
reduction predicates.  In (ii), minimality follows from the unit-$c_4$
argument above; in (i), the unit discriminant already attains the maximal
possible valuation among integral models.
\end{proof}

If $R\to S$ is a ring map and the maximal ideal $\mathfrak q$ of $S$
contracts to $\mathfrak p$ of $R$, then
\[
 f(r)\in\mathfrak q\iff r\in\mathfrak p.
\]
Hence the same signatures transfer through a chosen lying-over prime.  This
reduces the number-field completion theorem to the construction of the
integral ring map and its contraction identity.

\section{A uniform large-image theorem from Mazur and Tate inertia}

For a field $k$ put
\[
 U(t)=\begin{pmatrix}1&t\\0&1\end{pmatrix},
 \qquad
 L(t)=\begin{pmatrix}1&0\\t&1\end{pmatrix}.
\]

\begin{lemma}[Explicit unipotent factorization]
Let $A=\left(\begin{smallmatrix}a&b\\c&d\end{smallmatrix}\right)\in
\SL_2(k)$.
If $c\ne0$, then
\[
 A=U\!\left(\frac{a-1}{c}\right)L(c)
   U\!\left(\frac{d-1}{c}\right).
\]
If $c=0$, then $A$ is a product of four matrices of the form $U(t),L(t)$.
\end{lemma}

\begin{proof}
The first identity follows by direct multiplication and the equation
$ad-bc=1$.  In the triangular case, multiply by $L(1)$ to obtain a matrix with
nonzero lower-left entry, apply the first identity, and multiply back by
$L(-1)$.
\end{proof}

\begin{theorem}[Irreducible transvection criterion]
Let $\ell$ be prime and let $H\le\GL_2(\F_\ell)$.  If $H$ acts irreducibly on
$\F_\ell^2$ and contains a nontrivial transvection, then
\[
 \SL_2(\F_\ell)\le H.
\]
\end{theorem}

\begin{proof}
Choose a basis in which the transvection is $U(s)$ with $s\ne0$.  Its powers
contain every $U(t)$ because $\F_\ell$ is the prime field.  Irreducibility
supplies $A\in H$ moving the fixed line $\F_\ell e_1$, so the lower-left entry
$c$ of $A$ is nonzero.  Conjugating $A U(1)A^{-1}$ by
$U(-a/c)$ yields
\[
 L\!\left(-\frac{c^2}{\det A}\right),
\]
a nontrivial lower transvection.  Its powers give all $L(t)$.  The explicit
factorization now gives every determinant-one matrix.
\end{proof}

This argument avoids the classification of subgroups of $\GL_2(\F_\ell)$.
For a semistable elliptic curve over $\Q$, the required arithmetic inputs are:

\begin{enumerate}[label=(\arabic*)]
\item if $\ell>163$, the mod-$\ell$ representation is irreducible, because a
reducible representation gives a rational cyclic $\ell$-isogeny and Mazur's
prime-isogeny list is
\[
 \{2,3,5,7,11,13,17,19,37,43,67,163\};
\]
\item at a multiplicative prime $p\ne\ell$, Tate uniformization gives, in a
suitable basis,
\[
 \rho(\sigma)=
 \begin{pmatrix}1&N\,t_\ell(\sigma)\\0&1\end{pmatrix}
 \quad(\sigma\in I_p),
\]
where $N=\ord_p(q)=\ord_p(\Delta_{\min})$ and the tame character is
surjective modulo $\ell$.
\end{enumerate}
Thus $\ell\nmid N$ produces a nonzero transvection.

\section{A bounded Euclidean auxiliary prime}

\begin{theorem}[One-stage selector]
Let $B,N\in\Z_{>0}$ and put
\[
 A=B!N+1.
\]
Any prime divisor $\ell$ of $A$ satisfies
\[
 B<\ell,\qquad \ell\nmid N,\qquad \ell\le B!N+1.
\]
\end{theorem}

\begin{proof}
A prime $\ell\le B$ divides $B!$, hence cannot divide $B!N+1$.
Likewise, a prime dividing $N$ cannot divide $B!N+1$.  The upper bound follows
from $\ell\mid A$.
\end{proof}

The multiplicative inertia theorem also requires $\ell\ne p$.  Including $p$
in the avoidance product would introduce an unnecessarily large factor.  A
second Euclidean step avoids it without this loss.

\begin{theorem}[Two-stage selector]
Let $A_1=B!N+1$ and let $\ell_1$ be a prime divisor of $A_1$.  If
$\ell_1\ne p$, set $\ell=\ell_1$.  If $\ell_1=p$, choose a prime divisor
$\ell$ of
\[
 A_2=B!NA_1+1.
\]
Then
\[
 B<\ell,\qquad \ell\nmid N,
 \qquad \ell\ne p,
\]
and
\[
 \ell\le B!N(B!N+1)+1.
\]
\end{theorem}

\begin{proof}
Only the second case needs explanation.  The new prime divides neither $N$
nor $A_1$.  Since $p=\ell_1$ divides $A_1$, it cannot equal the new prime.
The displayed upper bound is the defining integer $A_2$.
\end{proof}

For fixed $B$,
\[
 B!N(B!N+1)+1
 \le \bigl((B!+1)(N+1)\bigr)^2,
\]
so
\[
 \log\ell\le
 2\log(B!+1)+2\log(N+1).
\]

\section{The local order is height-controlled}

Suppose $p^n\mid abc$.  Then
\[
 n\log p\le\log(abc).
\]
Since $a,b\le c$,
\[
 abc\le c^3,
 \qquad \log(abc)\le3\log c.
\]
Therefore
\[
 n\log2\le3h(a,b,c).
\]
At an odd multiplicative Frey prime,
\[
 N=\ord_p(\Delta_{\min})=2n,
\]
so
\[
 N\log2\le6h(a,b,c).
\]
Combining this with the Euclidean selector yields
\[
 \log\ell\le 2\log(1+h(a,b,c))+O_{\varepsilon}(1)
\]
after choosing
\[
 B\ge\max\left\{163,7,\left\lceil\frac{20}{\varepsilon}\right\rceil\right\}.
\]
The implied constant may depend on $\varepsilon$, as permitted in $abc$.

\section{Two exact absorption principles}

The verified public IUT IV specialization has the schematic form
\[
 h\le (1+\kappa)
       (\operatorname{Diff}+\operatorname{Cond}_{\mathrm{Frey}})
      +20\operatorname{Err}+O(1),
\]
where $\kappa=20d_{\mathrm{mod}}/\ell$.  For a rational Frey $j$-invariant,
$d_{\mathrm{mod}}=1$.

\begin{theorem}[Relative conductor absorption]
Assume that for every $\delta>0$ one can choose a genuine source such that
\[
 \kappa\le\delta
\]
and
\[
 (1+\delta)\operatorname{Diff}+20\operatorname{Err}
 \le \delta r+C_\delta
\]
uniformly in the $abc$ point.  Then the logarithmic $abc$ conjecture holds.
\end{theorem}

\begin{proof}
The Frey discriminant--conductor comparison gives
\[
 h\le(1+2\delta)r+O_\delta(1).
\]
Take $\delta=\varepsilon/2$.
\end{proof}

\begin{theorem}[Sublinear-height absorption]
Assume that for every $\delta,\eta>0$ one can choose a genuine source such
that $\kappa\le\delta$ and
\[
 (1+\delta)\operatorname{Diff}+20\operatorname{Err}
 \le \eta h+C_{\delta,\eta}.
\]
Then the logarithmic $abc$ conjecture holds.
\end{theorem}

\begin{proof}
Rearrangement gives
\[
 (1-\eta)h\le(1+\delta)r+O_{\delta,\eta}(1).
\]
For target $\varepsilon>0$, set
\[
 \delta=\frac{\varepsilon}{2},
 \qquad
 \eta=\frac{\varepsilon}{2(1+\varepsilon)}.
\]
Then
\[
 \frac{1+\delta}{1-\eta}=1+\varepsilon.
\]
\end{proof}

The elementary inequality
\[
 \log(1+x)\le \rho x+\rho-1-\log\rho
 \qquad(x\ge0,\ \rho>0)
\]
shows that every $O(\log(1+h))$ source term satisfies the second theorem.
Consequently the remaining quantitative target may be stated as:
\[
 (1+\delta)\operatorname{Diff}+20\operatorname{Err}
   =O_\delta(\log\ell+\log(N+1)+1).
\]

\section{The stacky sign quotient and the orbicurve square}

The object-level cover underlying the type $(1,\ell\text{-tors})$ construction
is
\[
 E\setminus E[\ell]\xrightarrow{[\ell]} E\setminus\{0\}.
\]
For every additive group $A$ one has the exact inverse-image identity
\[
 [n]^{-1}(A\setminus\{0\})=A\setminus A[n],
\]
and the map is equivariant for $x\mapsto-x$.

A coarse sign quotient is insufficient.

\begin{counterexample}[Failure of the coarse pullback]
In $A=\Z/6\Z$ with $n=3$, the distinct elements $1$ and $5=-1$ have the same
coarse sign orbit and both map to $3$.  Therefore
\[
 x\longmapsto([x],3x)
\]
is not injective, so the corresponding square of orbit sets is not
cartesian.
\end{counterexample}

The missing datum is the sign morphism.  For a group $G$ acting on $X,Y$ and
an equivariant map $f:X\to Y$, an object of the weak pullback of
$[X/G]\to[Y/G]$ along the atlas $Y\to[Y/G]$ is a quadruple
\[
 (x,y,g,\ g f(x)=y).
\]
Its normalized point is $gx$.  The assignment to $gx$ induces an equivalence
on isomorphism classes with $X$.  More strongly, a morphism between two such
objects exists exactly when their normalized points agree, and it is unique.
Thus the weak pullback groupoid is equivalent to the discrete groupoid on
$X$.

For the sign action, the stabilizer is trivial away from $A[2]$ and is the
full sign group at a two-torsion point.  This stacky inertia is precisely what
the coarse quotient erases.

The remaining geometric theorem is to realize the punctured torsion map as a
finite \emph{\'etale} map of algebraic curves, form the quotient stack, and
transport the resulting 2-cartesian square to the required Galois categories
and fundamental groups.

\section{The cyclic skeleton quotient}

An orientation-preserving deck transformation of an $\ell$-cycle is a
permutation of $\Z/\ell\Z$ commuting with $x\mapsto x+1$.  Such a permutation
is uniquely translation by its value at zero.  Hence
\[
 \operatorname{Deck}^+(C_\ell)\simeq\Z/\ell\Z,
\]
with canonical generator translation by one.  This is the finite graph
quotient expected from an $\ell$-fold cover of the oriented Tate/Berkovich
skeleton.

The remaining tempered theorem is an equivalence between the actual
rank-one quotient of the theta-root cover and this oriented deck group,
together with compatibility of the canonical graph cusp.

\section{Multiradial semantics and exact no-go results}

Let $\mathcal G$ denote genuine Theorem~3.11 outputs and $\mathcal S$ a
generated Ind1/Ind2/Ind3 syntax.  A semantic presentation consists of maps
\[
 \mathcal G\rightleftarrows\mathcal S
\]
whose represented regions include one another in the corresponding
directions.  Soundness and completeness identify the full possible-image
unions.  One-sided soundness already transfers the native lower branch.

\subsection{The distinguished branch}
For the lower inequality in Corollary~3.12, a classification of all genuine
outputs is unnecessary.  It suffices to produce one genuine output whose
represented region equals the native $q$-pilot region.  The complete candidate
Tate--Kummer model contains such an ordinary branch.  What remains is the
source theorem identifying this candidate branch with the actual
$\Theta^{\times\mu}_{\mathrm{LGP}}$ gluing output.

\subsection{Permutation quotient versus componentwise union}
Arbitrary label permutations preserve the total and uniform average of the
component volumes.  Therefore the average descends to the permutation-orbit
quotient.  They do not, however, preserve the volume of a componentwise union.

\begin{counterexample}[Orbit volume does not control union volume]
Take two labels and a two-point carrier.  Let the first component contain one
point and the second contain two.  The packet and its label swap both have
total cardinality $3$, whereas their componentwise union has total cardinality
$4$.
\end{counterexample}

Thus arbitrary Ind2 is valid only at a permutation-orbit/coarse numerical
level, unless a further correlated-packet theorem justifies a literal union.
The separate spectrum-preserving region model remains available for
componentwise inclusions.

\subsection{Scale rigidity}
If a nonempty set of positive $q$-log magnitudes is invariant under every
positive real scaling, then it contains every positive real number.  It is
therefore unbounded above and not bounded away from zero.  Hence an interface
that forgets scale up to unrestricted positive rescaling cannot support a
finite $q$-linked log-volume estimate.  The genuine log-Kummer theorem must
prove spectrum preservation or a specific finite envelope; unrestricted
scale-forgetting is excluded.

\section{Exact remaining theorem chain}

The reliable formal reductions leave the following source theorems.

\subsection*{Arithmetic global core}
\begin{enumerate}[label=G\arabic*.]
\item Construct a finite level-twelve torsion field for the Frey curve.
\item Either obtain $\sqrt{-1}$ by the Weil pairing or adjoin it explicitly;
prove the resulting finite extension and its Galois image.
\item Prove semistability after this extension, either by Raynaud's level
$3/4$ theorem or by direct ideal-signature transfer at all places.
\item Construct one odd multiplicative moduli place and identify its Tate
order with $2\ord_p(abc)$.
\end{enumerate}

\subsection*{Quantitative admissible prime}
\begin{enumerate}[label=P\arabic*.]
\item Formalize reducibility $\Rightarrow$ rational cyclic isogeny and Mazur's
prime-isogeny classification.
\item Formalize the Tate-inertia matrix and surjective tame character.
\item Transport the actual $\GL_2$ representation image to the explicit
matrix model.
\item Prove solvability of the fixed level-twelve (and optional quadratic)
extension and apply perfect-image restriction.
\item Prove a logarithmic or relative bound for the genuine IUT
$\operatorname{Diff}/\operatorname{Err}$ terms.
\end{enumerate}

\subsection*{Anabelian and tempered fiber}
\begin{enumerate}[label=A\arabic*.]
\item Realize the punctured torsion map as a finite \'etale cover.
\item Construct the sign quotient stack/orbicurve and its 2-cartesian square.
\item Identify the corresponding \'etale fundamental groups, core, and type
predicates.
\item Construct the local Tate/Berkovich skeleton and theta-root cover.
\item Identify the rank-one tempered quotient and canonical graph cusp.
\end{enumerate}

\subsection*{Genuine multiradial source}
\begin{enumerate}[label=M\arabic*.]
\item Construct actual Hodge-theater/log-Kummer output objects.
\item Prove the distinguished gluing branch equals the native $q$-pilot.
\item Prove a sound semantic presentation by the generated syntax.
\item Prove Ind1/Ind2 volume invariance in the correct quotient semantics.
\item Prove Ind3 upper semi-compatibility and the all-output envelope.
\item Prove native procession-volume calibration and the public IUT IV
component formula.
\end{enumerate}

Only after G1--G4, P1--P5, A1--A5, and M1--M6 are established independently
and connected to the already verified downstream Lean chain is it legitimate
to conclude an unconditional proof of the $abc$ conjecture.

\section*{Lean correspondence}
The principal modules associated with the results of this manuscript include:
\begin{itemize}
\item \texttt{FreyOddPrimeSupport},
  \texttt{MinimalModelFromUnitC4},
  \texttt{LocalIntegralModelReduction},
  \texttt{LocalReductionLiesOverTransfer};
\item \texttt{SL2UnipotentFactorization},
  \texttt{SL2IrreducibleTransvectionCriterion},
  \texttt{TateInertiaTransvectionReduction},
  \texttt{MazurPrimeIsogenyBoundary};
\item \texttt{EuclidAuxiliaryPrimeSelector},
  \texttt{EuclidAuxiliaryPrimeAvoidOne},
  \texttt{EuclidAuxiliaryPrimeLogBound},
  \texttt{FreyLocalOrderHeightBound};
\item \texttt{RelativeDifferentErrorAbsorption},
  \texttt{SublinearHeightSourceAbsorption},
  \texttt{LogOnePlusSublinear};
\item \texttt{PuncturedTorsionCover},
  \texttt{OrbifoldQuotientNecessity},
  \texttt{ActionGroupoidPullbackClasses},
  \texttt{ActionGroupoidPullbackMorphisms},
  \texttt{SignStabilizerInertia};
\item \texttt{TemperedCycleSkeleton},
  \texttt{MultiradialPresentationTransfer},
  \texttt{DistinguishedMultiradialBranch},
  \texttt{Ind2PermutationVolumeQuotient}, and
  \texttt{ValueGroupScaleRigidity}.
\end{itemize}
Each module is merged into the main branch only after the Lean kernel build
succeeds.

\end{document}
